\begingroup\small
\begin{thebibliography}{99}

\bibitem{Salehi2002}
S. Salehi,
\emph{Decidable Formulas of Intuitionistic Primitive Recursive Arithmetic},
Reports on Mathematical Logic \textbf{36} (2002), 55--61.
Fact 3.5, p.~57, states the arbitrary-matrix primitive recursive
witness theorem used here.
\url{https://rml.tcs.uj.edu.pl/rml-36/36-SALEHI.pdf}.

\bibitem{Wehmeier1997}
K. F. Wehmeier,
\emph{Fragments of HA based on $\Sigma_1$-induction},
Archive for Mathematical Logic \textbf{37} (1997), 37--49.
\url{https://doi.org/10.1007/s001530050081}.
The publisher abstract and metadata were checked; the witness formulation
used in this report is the explicit statement in Salehi.

\bibitem{CorvajaLombardoZannier2024}
P. Corvaja, D. Lombardo, and U. Zannier,
\emph{Examples of effectivity for integral points on certain curves of genus 2},
arXiv:2411.17930v2, revised 28 December 2024.
Introduction and Remark 1.2.
\url{https://arxiv.org/abs/2411.17930};
\url{https://arxiv.org/pdf/2411.17930}.

\bibitem{BiluIllengo2009}
Y. Bilu and M. Illengo,
\emph{Effective Siegel's theorem for modular curves},
arXiv:0905.0418 (2009), Introduction and Theorem 1.2.
\url{https://arxiv.org/pdf/0905.0418}.

\bibitem{CorvajaZannier2002}
P. Corvaja and U. Zannier,
\emph{A subspace theorem approach to integral points on curves},
Comptes Rendus de l'Acad\'emie des Sciences, Paris, S\'erie I
\textbf{334} (2002), 267--271.
\url{https://doi.org/10.1016/S1631-073X(02)02240-9}.
Open text:
\url{https://www.numdam.org/item/10.1016/S1631-073X(02)02240-9.pdf}.

\bibitem{Evertse2010}
J.-H. Evertse,
\emph{On the quantitative Subspace Theorem},
arXiv:1008.2268 (2010), especially Theorems 2.1--2.2.
\url{https://arxiv.org/pdf/1008.2268}.
Used to identify the quantitative dependency, not as an iPRA formalization.

\bibitem{Levin2016}
A. Levin,
\emph{Integral points and the Shafarevich conjecture for abelian surfaces},
in \emph{Diophantische Approximationen},
Oberwolfach Report 21/2016, 1176--1178, especially 1176--1177.
\url{https://doi.org/10.4171/OWR/2016/21}.
This source further explains the effectivity distinction for genus two
with different choices of poles.

\bibitem{ProveIt2026}
V. Reshetnikov and contributors,
\emph{ProveIt}, repository snapshot:
\begin{quote}\small\ttfamily
7eeb3c0e821594d5d5ebcfbaf4bdba64819e3502.
\end{quote}
\url{https://github.com/VladimirReshetnikov/ProveIt}.
Inspected foundation files:
\begin{quote}\small\ttfamily
Logic/PrimitiveRecursiveArithmetic/ConciseIPRA/Lean/\\
ConciseIPRA/PR.lean\\
ConciseIPRA/Deduction.lean\\
ConciseIPRA/NativeBiInterpretation.lean
\end{quote}
Related packet index and status files are recorded with their exact
paths and hashes in \texttt{provenance/repository\_manifest.json}.
No upstream modification or full upstream build is asserted.

\bibitem{UserReference}
\emph{Siegel's theorem on integral points}, Wikipedia,
the reference supplied with the request; accessed 15 November 2026.
\url{https://en.wikipedia.org/wiki/Siegel\%27s_theorem_on_integral_points}.
The primary sources above support the mathematical and logical
status discussion.

\end{thebibliography}
\endgroup

